\exercisesection{Exercises for § 5}

\begin{enumerate}
\def\labelenumi{\arabic{enumi})}
\item
  Determine the algebra of symmetric invariants of a finite dihedral group (for its canonical representation of dimension 2, cf.~§4, no. 2).
\item
  Let A be a principal ring, E a free A-module of finite rank l, and G a finite subgroup of \(\mathbf{GL}(\mathbf{E})\) . Assume the following:
\end{enumerate}

\begin{enumerate}
\def\labelenumi{(\roman{enumi})}
\item
  If \(q = \operatorname{Card}(G)\) , the element \(q.1\) of \(A\) is invertible.
\item
  G is generated by pseudo-reflections in E (i.e.~by elements s such that \((s - 1)(E)\) is a monogenic submodule of E, cf.~§2, Exerc. 1.)
\end{enumerate}

Let \(\mathbf{S}(\mathrm{E})\) be the symmetric algebra of \(\mathrm{E}\) , and let \(\mathbf{S}(\mathrm{E})^{\mathrm{G}}\) be the subalgebra of \(\mathbf{S}(\mathrm{E})\) consisting of the elements invariant under \(\mathrm{G}\) . Show that, for any homomorphism from \(\mathrm{A}\) into a field \(k\) , \(\mathbf{S}(\mathrm{E})^{\mathrm{G}} \otimes k\) can be identified with \(\mathbf{S}(\mathrm{E} \otimes k)^{\mathrm{G}}\) . Deduce, by applying Th. 4, that \(\mathbf{S}(\mathrm{E})^{\mathrm{G}}\) is a graded polynomial algebra over \(\mathrm{A}\) .

¶ 3) Let K be a field of characteristic zero, V a vector space of finite dimension l over K, and G a subgroup of GL(V) generated by pseudo-reflections. Put q = Card(G); denote by S (resp. L) the symmetric algebra (resp. exterior algebra) of V. If \(x \in V\) , denote by x (resp. \(x'\) ) its canonical image in S (resp. L).

\begin{enumerate}
\def\labelenumi{\alph{enumi})}
\item
  Let \(E = S \otimes L\) be the tensor product of S and L. Show that there exists a unique derivation d on E such that \(dx = x'\) and \(dx' = 0\) for all \(x \in V\) .
\item
  Let \(S^{G}\) be the algebra of symmetric invariants of G and let \(P_{1},\ldots,P_{l}\) be homogeneous elements of \(S^{G}\) such that \(S^{G}=K[P_{1},\ldots,P_{l}]\) . For every subset \(I=\{i_{1},\ldots,i_{r}\}\) of \([1,l]\) with \(i_{1}<\cdots<i_{r}\) , put
\end{enumerate}

\[
\omega_ {\mathrm{I}} = d \mathrm{P} _ {i _ {1}} \dots d \mathrm{P} _ {i _ {r}}.
\]

Show that the \(\omega_{I}\) are linearly independent over S, and that they belong to the subalgebra \(E^{G}\) consisting of the elements of E invariant under G. Deduce that, for all \(\omega \in E^{G}\) , there exist \(a, c_{I} \in S^{G}\) , with \(a \neq 0\) , such that

\[
a \omega = \sum_ {\mathrm{I}} c _ {\mathrm{I}} \omega_ {\mathrm{I}}.
\]

\(c)\) Show that every element of \(\mathrm{E}^{\mathrm{G}}\) contained in \(S\otimes \bigwedge^{l}\mathrm{V}\) is of the form \(c.dP_{1}\ldots dP_{l}\) , with \(c\in S^{\mathrm{G}}\) . (Apply Prop. 5.)

\begin{enumerate}
\def\labelenumi{\alph{enumi})}
\setcounter{enumi}{3}
\item
  Show that the \(\omega_{I}\) form a basis of the \(S^{G}\) -module \(E^{G}\) . (With the notation in b), multiply both sides of the relation \(a\omega = \sum_{I} c_{I}\omega_{I}\) by an element \(\omega_{J}\) , and apply c); deduce that, if I is the complement of J, a divides \(c_{I}\) ; hence \(^{9}\) the fact that the \(\omega_{I}\) generate \(E^{G}\) .)
\item
  Let \(S_{n}\) (resp. \(L_{m}\) ) be the homogeneous component of \(S\) (resp. \(L\) ) of degree \(n\) (resp. \(m\) ). Put
\end{enumerate}

\[
\begin{array}{r l r} & & {\mathrm{E} _ {n, m} = \mathrm{S} _ {n} \otimes \mathrm{L} _ {m}, \quad \mathrm{E} _ {n, m} ^ {\mathrm{G}} = \mathrm{E} ^ {\mathrm{G}} \cap \mathrm{E} _ {n, m}, \quad a _ {n, m} = \dim \mathrm{E} _ {n, m} ^ {\mathrm{G}},} \\ & & {a (\mathrm{X}, \mathrm{Y}) = \sum_ {n, m \geqslant 0} a _ {n, m} \mathrm{X} ^ {n} \mathrm{Y} ^ {m}.} \end{array}
\]

By using \(d\) ), prove the formula

\[
a (\mathrm{X}, \mathrm{Y}) = \prod_ {i = 1} ^ {l} \frac {1 + \mathrm{Y} . \mathrm{X} ^ {p _ {i} - 1}}{1 - \mathrm{X} ^ {p _ {i}}},
\]

where \(p_{i} = \deg(P_{i})\) .

\(f)\) If \(g \in \mathbf{G}\) , let \(\operatorname{Tr}_{n,m}(g)\) be the trace of the automorphism of \(\mathbf{E}_{n,m}\) defined by \(g\) ; put

\[
\operatorname{Tr} (\mathrm{X}, \mathrm{Y}) (g) = \sum_ {n, m \geqslant 0} \operatorname{Tr} _ {n, m} (g) \mathrm{X} ^ {n} \mathrm{Y} ^ {m}.
\]

Show that

\[
\operatorname{Tr} (X, Y) (g) = \frac {\det (1 + Y g)}{\det (1 - X g)}.
\]

\begin{enumerate}
\def\labelenumi{\alph{enumi})}
\setcounter{enumi}{6}
\tightlist
\item
  Let \(q = \operatorname{Card}(\mathbf{G})\) . Show that
\end{enumerate}

\[
\frac {1}{q} \sum_ {g \in G} \operatorname{Tr} (X, Y) (g) = a (X, Y),
\]

i.e.

\[
\frac {1}{q} \sum_ {g \in G} \frac {\det (1 + Y g)}{\det (1 - X g)} = \prod_ {i = 1} ^ {l} \frac {1 + Y . X ^ {p _ {i} - 1}}{1 - X ^ {p _ {i}}}.
\]

(Use the following result: if \(\mathbf{G}\) acts on a finite dimensional vector space \(\mathbf{E}\) , the dimension of the space of elements of \(\mathbf{E}\) invariant under \(\mathbf{G}\) is equal to \(\frac{1}{q}\sum_{g\in \mathbf{G}}\mathrm{Tr}_{\mathbf{E}}(g).\) )

What does this formula give for Y = 0?

\begin{enumerate}
\def\labelenumi{\alph{enumi})}
\setcounter{enumi}{7}
\tightlist
\item
  For any integer \(p \geqslant 0\) , let \(H_{p}\) be the set of elements \(g \in G\) that have 1 as an eigenvalue with multiplicity p.~Let \(h_{p} = \text{Card}(H_{p})\) . Prove the formula
\end{enumerate}

\[
\sum_ {p = 0} ^ {l} h _ {p} \mathbf {T} ^ {p} = \prod_ {i = 1} ^ {l} (p _ {i} - 1 + \mathbf {T}).
\]

(In the formula in \(g\) ) above, replace \(Y\) by \(-1 + T(1 - X)\) and then put \(X = 1\) in the result. If \(g \in H_p\) , the term \(\operatorname{Tr}(X, Y)(g)\) becomes \(T^p\) .)

¶ 4) *Let \(G _ {1} = \mathbf {S L} _ {3} (\mathbf {F} _ {2})\).

\begin{enumerate}
\def\labelenumi{\alph{enumi})}
\item
  Show that \(\mathbf{G}_1\) is a non-abelian simple group of order 168, containing 21 elements of order 2.
\item
  Show that the degrees of the irreducible complex representations of \(G_{1}\) are 1, 3, 3, 6, 7, 8.
\item
  Let \(\rho: \mathbf{G}_1 \to \mathbf{GL}_3(\mathbf{C})\) be an irreducible representation \(^{10}\) of \(\mathbf{G}_1\) of degree 3. If \(y \in \mathbf{G}_1\) is of order 2, show that \(\operatorname{Tr}(\rho(y)) = -1\) . Deduce that \(-\rho(y)\) is a reflection.
\item
  Let G be the subgroup of \(\mathbf{GL}_{3}(\mathbf{C})\) generated by the elements \(-\rho(y)\) , for y of order 2 in \(G_{1}\) . Show that G is isomorphic to \(G_{1} \times \{1, -1\}\) , and is thus of order 336.
\item
  Show that the characteristic degrees \(k_{1}, k_{2}, k_{3}\) of the algebra of symmetric invariants of \(\mathbf{G}\) are equal to 4, 6 and 14. (Use the relations \(\prod_{i} k_{i} = 336\) and \(\sum_{i} (k_{i} - 1) = 21\) .)
\item
  Show that G is not a Coxeter group.*
\end{enumerate}

\begin{enumerate}
\def\labelenumi{\arabic{enumi})}
\setcounter{enumi}{4}
\tightlist
\item
  *Let \(K\) be a field and \(S = K[X_1, \ldots, X_n]\) a graded polynomial algebra over \(K\) , generated by algebraically independent elements \(X_i\) , homogeneous of degrees \(> 0\) .
\end{enumerate}

\begin{enumerate}
\def\labelenumi{\alph{enumi})}
\tightlist
\item
  Let \(Y_{1},\ldots,Y_{n}\) be homogeneous elements of S of degrees \textgreater0, and let \(R = K[Y_{1},\ldots,Y_{n}]\) be the subalgebra of S generated by these elements. Prove the equivalence of the following properties:
\end{enumerate}

\begin{enumerate}
\def\labelenumi{(\roman{enumi})}
\item
  \((\mathrm{Y}_1,\dots ,\mathrm{Y}_n)\) is an S-regular sequence.
\item
  S is integral over R.
\item
  The ideal of \(S\) generated by \(Y_{1},\ldots ,Y_{n}\) is of finite codimension in \(S\) .
\item
  For any extension \(\overline{\mathbf{K}}\) of \(\mathbf{K}\) , the system of equations \(Y_{i}(x_{1},\ldots ,x_{n}) = 0\) ( \(1\leqslant i\leqslant n\) , \(x_{i}\in \overline{\mathbf{K}}\) ) has only the trivial solution \((0,\dots ,0)\) .
\end{enumerate}

\(((\mathrm{i}) \Longleftrightarrow (\mathrm{iii})\) follows from a theorem of Macaulay (Commutative Algebra); \((\mathrm{iii}) \Longleftrightarrow (\mathrm{iv})\) follows from the Zeros Theorem (Commutative Algebra, Chap. V, §3, no. 3, Prop. 2); (ii) \(\Longleftrightarrow (\mathrm{iii})\) is easy.)

If these properties are satisfied, show that the \(Y_{i}\) are algebraically independent over K, and that S is a free R-module of rank equal to

\[
\frac {\prod_ {i} \deg (\mathrm{Y} _ {i})}{\prod_ {i} \deg (\mathrm{X} _ {i})}.
\]

\begin{enumerate}
\def\labelenumi{\alph{enumi})}
\setcounter{enumi}{1}
\tightlist
\item
  Let G be a finite group of automorphisms of the graded algebra S, let \(S^{G}\) be its subalgebra of invariants, and let \(Y_{1},\ldots,Y_{n}\) be elements of \(S^{G}\) satisfying conditions (i),\ldots, (iv) above. Show that \(S^{G}=K[Y_{1},\ldots,Y_{n}]\) if and only if
\end{enumerate}

\[
\operatorname{Card} (\mathrm{G}) = \frac {\prod_ {i} \deg \left(\mathrm{Y} _ {i}\right)}{\prod_ {i} \deg \left(\mathrm{X} _ {i}\right)}. *
\]

¶ 6) *Let n be an integer ≥ 1, q a power of a prime number, K = F\_q, V = K\^{}n, G = GL(n, K) and G\_1 = SL(n, K). Identify G with the group GL(V). Further, denote by S the algebra S(V) = K{[}X\_1, \ldots, X\_n{]} and R (resp. R\_1) the subalgebra of S consisting of the elements invariant under G (resp. G\_1).

\begin{enumerate}
\def\labelenumi{\alph{enumi})}
\tightlist
\item
  Let \(\mathbf{e} = (e_1, \ldots, e_n)\) be a sequence of integers \(\geqslant 0\) . Put
\end{enumerate}

\[
\mathrm{L} _ {\mathbf {e}} = \det (\mathrm{X} _ {i} ^ {q ^ {e _ {j}}}).
\]

This is an element of \(S\) .

Show that

\[
g. \mathrm{L} _ {\mathbf {e}} = \det (g) \mathrm{L} _ {\mathbf {e}} \text {for all} g \in \mathrm{G}.
\]

(Remark that, if \(g.X_i = \sum_j a_{ij} X_j\) , then \(g.X_i^{q^e} = \sum_j a_{ij} X_j^{q^e}\) .) In particular, \(L_e\) belongs to \(R_1\) .

\begin{enumerate}
\def\labelenumi{\alph{enumi})}
\setcounter{enumi}{1}
\tightlist
\item
  Let \(j \in [1, n]\) . Put
\end{enumerate}

\[
Z _ {j} = \prod_ {(a _ {i j})} (X _ {j} + \sum_ {i > j} a _ {i j} X _ {i}),
\]

the product being over all families \((a_{ij})_{j < i\leqslant n}\) of elements of K. Let

\[
\mathrm{T} = \prod_ {j = 1} ^ {n} Z _ {j}.
\]

Show that T divides all the \(L_{e}\) . Deduce that \(T = L_{e_{n}}\) , where \(e_{n} = (0, 1, \ldots, n - 1)\) , and that T is invariant under \(G_{1}\) .

\begin{enumerate}
\def\labelenumi{\alph{enumi})}
\setcounter{enumi}{2}
\tightlist
\item
  Denote by \(\mathbf{e}_i\) ( \(1 \leqslant i \leqslant n - 1\) ) the sequence \((0, 1, \ldots, i - 1, i + 1, \ldots, n)\) , and put \(\mathrm{Y}_i = \mathrm{L}_{\mathbf{e}_i} / \mathrm{T}\) , cf.~b). Show that the \(\mathrm{Y}_i\) belong to R and that \(\deg(\mathrm{Y}_i) = q^n - q^i\) .
\end{enumerate}

\(d)\) Let \(S' = K[X_1, \ldots, X_{n-1}]\) , and let \(T', Y_1', \ldots, Y_{n-2}'\) be the elements of \(S'\) defined in the same way as \(T, Y_1, \ldots, Y_{n-1}\) (but replacing \(n\) by \(n-1\) ). Let \(f: S \to S'\) be the homomorphism defined by

\[
f \left(\mathrm{X} _ {i}\right) = \mathrm{X} _ {i} (1 \leqslant i \leqslant n - 1), f \left(\mathrm{X} _ {n}\right) = 0.
\]

Show that

\[
f (\mathrm{T}) = 0, \quad f (\mathrm{Y} _ {1}) = \mathrm{T} ^ {\prime q (q - 1)}, \quad f (\mathrm{Y} _ {i}) = \mathrm{Y} _ {i - 1} ^ {\prime q} \quad \text{for} \quad 2 \leqslant i \leqslant n - 1.
\]

\begin{enumerate}
\def\labelenumi{\alph{enumi})}
\setcounter{enumi}{4}
\tightlist
\item
  Show that the family \((T, Y_{1}, \ldots, Y_{n-1})\) satisfies condition (iv) of Exerc. 5. (Let \(x = (x_{1}, \ldots, x_{n})\) be a zero of the system \((T, Y_{1}, \ldots, Y_{n-1})\) in an extension \(\overline{K}\) of K. Since x is a zero of T, the \(x_{i}\) satisfy at least one non-trivial linear relation with coefficients in K. Transforming x by an element of \(G_{1}\) if necessary, we can assume that \(x_{n} = 0\) . Conclude by applying d) and arguing by induction on n.)
\end{enumerate}

\(f)\) Show that \(\mathbf{R}_1 = \mathbf{K}[\mathbf{T},\mathbf{Y}_1,\dots ,\mathbf{Y}_{n - 1}]\) , and that \(\mathrm{T},\mathrm{Y}_1,\dots ,\mathrm{Y}_{n - 1}\) are algebraically independent (``Dickson's Theorem'' - apply e) and Exerc. 5, remarking that the order of \(\mathbf{G}_1\) is equal to \(\deg (\mathrm{T}) = \sum_{i}\deg (\mathrm{Y}_i)\)).

\begin{enumerate}
\def\labelenumi{\alph{enumi})}
\setcounter{enumi}{6}
\tightlist
\item
  Show that \(R = K[T^{q-1}, Y_{1}, \ldots, Y_{n-1}]\) .
\end{enumerate}

¶ 7) *Let R be a regular local ring (Commutative Algebra, Chap. VIII, §5, no. 1), with maximal ideal m and residue field k. Let G be a finite group of automorphisms of R, and let \(R' = R^{G}\) be the subring of R consisting of the elements invariant under G. This is a local ring with maximal ideal \(m' = m \cap R'\) . Assume the following:

\begin{enumerate}
\def\labelenumi{(\roman{enumi})}
\item
  \(R'\) is noetherian and R is an \(R'\) -module of finite type.
\item
  The composite \(R'\to R\to k\) is surjective.
\end{enumerate}

Put \(V = \mathfrak{m} / \mathfrak{m}^2\) ; this is a \(k\) -vector space. The action of \(G\) on \(R\) defines a homomorphism \(\varepsilon : G \to \mathbf{GL}(V)\) .

\begin{enumerate}
\def\labelenumi{\alph{enumi})}
\item
  Let \(\mathfrak{p}\) be a prime ideal of \(R\) of height 1 (Commutative Algebra, Chap. VII, §1, no. 6) and let \(s \in G\) be such that \(s(\mathfrak{p}) = \mathfrak{p}\) and that \(s\) acts trivially on \(R / \mathfrak{p}\) . Show that \(\varepsilon(s)\) is a pseudo-reflection of \(V\) . (Remark that the image of \(\mathfrak{p}\) in \(\mathfrak{m} / \mathfrak{m}^2\) is of dimension 0 or 1.)
\item
  Show that, if \(R'\) is regular, the subgroup \(\varepsilon(G)\) of \(\mathbf{GL}(V)\) is generated by pseudo-reflections, (Let H be the subgroup of G generated by elements whose image under \(\varepsilon\) is a pseudo-reflection, and let \(R^{H}\) be the subring of R consisting of elements invariant under H. Show, using a), that no prime ideal of \(R'\) of height 1 is ramified in \(R^{H}\) ; using the fact that \(R^{H}\) is integrally closed, deduce by means of the Purity Theorem \(^{11}\) that \(R' = R^{H}\) , so H = G.)
\item
  Assume now that the order of \(\mathbf{G}\) is coprime to the characteristic of \(k\) . Show that \(\varepsilon\) is injective.
\end{enumerate}

Let \((\mathfrak{m}_{n}^{\prime})\) be the filtration of \(R^{\prime}\) induced by the m-adic filtration \((\mathfrak{m}^{n})\) of R, and let

\[
i: \mathrm{gr} (\mathrm{R} ^ {\prime}) \to \mathrm{gr} (\mathrm{R})
\]

be the canonical homomorphism of the associated graded algebra of \(R'\) to that of R (Commutative Algebra, Chap. III, §2). Show that i is injective, and that its image is the subring \(\mathrm{gr}(R)^{\mathrm{G}}\) of \(\mathrm{gr}(R)\) consisting of the elements invariant under G.

\(d)\) Retain the assumptions and notation of \(c\) ), and assume further that \(\varepsilon (\mathbf{G})\) is generated by pseudo-reflections. If \(l = \dim V\) , let \(\mathrm{P}_1,\ldots ,\mathrm{P}_l\) be algebraically independent homogeneous generators of the \(k\) -algebra \(\mathrm{gr}(R)^{\mathrm{G}}\) (such elements exist by Th. 4 and the fact that \(\mathrm{gr}(R)\) can be identified with the symmetric algebra of V); let \(p_1,\ldots ,p_l\) be their degrees. By \(c\) ), we can find \(x_{i}\in \mathfrak{m}_{p_{i}}^{\prime}\) such that \(\mathrm{gr}(x_i) = \mathrm{P}_i\) . Show that the \(x_{i}\) generate the ideal \(\mathfrak{m}'\) and deduce that \(R'\) is regular.*

¶ 8) *Let V be a finite dimensional vector space over a field K, let S be the symmetric algebra of V, and let G be a finite subgroup of GL(V). Assume that the algebra \(S^G\) of symmetric invariants of \(G\) is a graded polynomial algebra.

\begin{enumerate}
\def\labelenumi{\alph{enumi})}
\tightlist
\item
  Let u be an element of the dual \(V^{*}\) of V, and let \(G_{u}\) be the subgroup of G consisting of the elements leaving u invariant. Show that \(G_{u}\) is generated by pseudo-reflections. (The linear form u extends to a homomorphism \(f_{u}: S \to K\) . If \(S_{u}\) denotes the local ring of S with respect to the kernel of \(f_{u}\) , the ring \(S_{u}\) is regular. Conclude by applying Exerc. 7 b) to \(G_{u}\) , considered as a group of automorphisms of \(S_{u}\) .)
\end{enumerate}

In particular, G is generated by pseudo-reflections.

\begin{enumerate}
\def\labelenumi{\alph{enumi})}
\setcounter{enumi}{1}
\tightlist
\item
  Let A be a subset of \(V^{*}\) , and let \(G_{A} = \bigcap_{u \in A} G_{u}\) . Show that \(G_{A}\) is generated by pseudo-reflections. (Extending the base field if necessary, we can assume that K is infinite. Show that in that case there exists an element v of the vector subspace of \(V^{*}\) generated by A such that \(G_{v} = G_{A}\) . Conclude by applying a) to \(v\) .)
\end{enumerate}

\begin{enumerate}
\def\labelenumi{\arabic{enumi})}
\setcounter{enumi}{8}
\tightlist
\item
  Let V be a vector space of dimension 4 over a finite field K, of characteristic different from 2, and let Q be a non-degenerate quadratic form on V of index 2 (Algebra, Chap. IX, §4, no. 2). Let G = O(Q) be the orthogonal group of this form; this is a finite group generated by reflections (loc. cit., §6, no. 4, Prop. 5).
\end{enumerate}

\begin{enumerate}
\def\labelenumi{\alph{enumi})}
\item
  Let E be a maximal totally isotropic subspace of V, and let \(G_{E}\) be the subgroup of G consisting of the elements g such that \(g(x) = x\) for all \(x \in E\) . Show that \(G_{E}\) is isomorphic to the additive group of K, and contains no pseudo-reflection.
\item
  Show that the algebra of symmetric invariants of G is not a graded polynomial algebra (use the preceding exercise).
\end{enumerate}

